\chapter{鲁棒性}
\label{chap:trust-robustness}

邮件分类器在旧邮件上表现良好，换到新月份却频繁误拦。若它还读取扫描附件，图片压缩可能使文字识别出错，改变模型读到的内容；正常邮件也可能开始使用过去常见于垃圾邮件的模板，使旧分类规则失效。前一种变化发生在信息采集环节，后一种变化改变了特征与标签的关系。只有知道需要承受哪一种变化，才能判断模型是否鲁棒。

\term{鲁棒性}（robustness）描述模型在规定变化范围内保持有效表现的能力。它必须附带变化对象、范围和评价标准。系统还需要在证据不足或模型失效时识别变化、拒绝高风险输入并回退；本章把这种能力称为\term{变化应对能力}（resilience to change）。前者回答模型在变化后是否仍正确，后者回答系统在不能确认模型正确时怎样限制损害，二者必须分别报告。评估一种变化后，既要检查预测性能损失了多少，也要区分能否证明局部稳定、能否识别已经超出适用范围的输入。改进方法则应针对诊断出的薄弱环节：改变训练约束、利用新环境信息，或限制缺乏证据的自动处理。分布变化下的可学习性条件见第~\ref{chap:transfer-learning-theory}章；这里重点考察这些条件如何影响模型的检验与改进。

沿用基础章的三个问题，先规定模型应当承受哪些变化，以及变化后至少保持怎样的表现。再通过新环境中的风险测试、规定邻域内的数学认证和拒绝效果评价，检查这些要求是否有依据。若数据来源或运行中的模型继续变化，还要监测原结论是否失效，并准备回退和重新验证。

最不利扰动既可以用于描述测量误差，也可以由攻击者主动寻找；本章围绕变化范围与模型性能讨论共同基础，第~\ref{chap:trust-security}章再加入攻击权限、资源和损害路径。这样的区分使局部稳定性方法能够共用，却不把全部安全问题缩减为输入扰动。

\section{模型的鲁棒性}
\label{sec:robustness-fragility}

变化本身不是脆弱性：图像压缩后仍能正确分类，就没有在该变化下暴露失效。本节分别确定“要求模型承受什么变化”和“模型为何未能承受变化”。前者给出评价范围，后者解释性能下降的机制；同一种变化可以暴露不同机制，同一种捷径也可能被多种变化触发。

\subsection{变化范围与性能要求}
\label{sec:robustness-change-scope}

设输入$\bm{x}$的允许变化构成集合$U(\bm{x})$。如果这些变化保留标签$y$，可以用最不利变化下的损失评价模型：
\begin{equation}
 R_{\mathrm{pert}}(f)
 =\mathbb E_{(\bm{x},y)\sim\mathcal P}
   \left[\sup_{\bm{x}'\in U(\bm{x})}
   \ell(f(\bm{x}'),y)\right].
 \label{eq:robustness-perturbation-risk}
\end{equation}
集合$U(\bm{x})$可以描述像素扰动、文本中的等义修改或测量误差，但它的合理性需要由任务确定。删除一句否定词可能只改变很少字符，却改变正确标签；此时要求输出不变本身就是错误的鲁棒性目标。

局部扰动并未涵盖全部分布变化。输入比例改变、标签比例改变和条件标签规律改变具有不同的信息要求，具体分类见第~\ref{sec:distribution-shift}节。对于回归或分类器$f$，环境$e$中的风险记为$R_e(f)$。鲁棒性可以要求所有允许环境的风险均不超过容限，也可以要求相对于原环境的下降不超过给定幅度。绝对风险和下降幅度需要一起看：一个始终表现很差的模型，也可能几乎没有下降。

\begin{example}[总体表现掩盖群体脆弱性]
\label{ex:robustness-mixture}
在构造的旧环境中，99\%的邮件属于常见模板，其错误率为1\%；其余1\%属于少见模板，错误率为50\%。总体错误率为
$0.99\times0.01+0.01\times0.5=1.49\%$。
若新环境中少见模板的比例增至30\%，而各组内部错误率不变，总体错误率就变为15.7\%。失效并不要求模型参数或组内标签规律发生变化。
\end{example}

图~\ref{fig:trust-robust-neighborhood}区分两类证据：找到一个允许范围内的错误输入，就能否定相应鲁棒性主张；确认整个邻域都稳定，则需要覆盖该范围的认证依据。后面的检验方法会分别处理这两个问题。

\begin{figure}[!htb]
  \centering
  % 集合关系教学示意：外圈允许变化，内圈已有认证，方点反例。
\begin{tikzpicture}[figure picture,foundation picture]
 \path[use as bounding box] (0,-1) rectangle (110,58);
 \fill[FigureBlueFill] (3,2) rectangle (58,51);
 \fill[FigureRedFill] (58,2) rectangle (107,51);
 \draw[draw=FigureStroke,line width=1.1pt] (58,2)--(58,51);
 \node at (25,55) {$\hat y=0$};
 \node at (83,55) {$\hat y=1$};
 \draw[draw=FigureBlue,line width=1pt,dashed] (39,26) circle (23);
 \filldraw[draw=FigureYellow,fill=FigureYellowFill,line width=1pt]
  (39,26) circle (11);
 \fill[FigureInk] (39,26) circle (1.5);
 \node[anchor=north] at (39,23) {$\bm x$};
 \draw[foundation flow,draw=FigureRed] (42,26)--(59,26);
 \draw[draw=FigureRed,fill=white,line width=1pt] (59.5,24.5) rectangle (62.5,27.5);
 \node[anchor=south west] at (63,28) {$\bm x'$};
 \node[anchor=east] at (15,45) {$U(\bm x)$};
\end{tikzpicture}

  \caption{允许变化、认证范围与搜索反例的集合关系示意。虚线圈表示$U(\bm x)$，内圈表示已有认证支持的稳定范围，方点$\bm x^{\prime}$是跨越决策边界的允许变化。圆的尺寸不是随机平滑的实验结果；搜索没有找到反例时，也不能直接把整个允许集合视为已认证。}
  \label{fig:trust-robust-neighborhood}
\end{figure}

\subsection{失效机制与信息缺口}
\label{sec:robustness-failure-mechanisms}

给出变化范围后，还不能断定模型会失败。局部扰动可能跨越距离输入很近的决策边界；群体比例变化可能放大原来就存在的高错误率；特征与标签关系变化则可能使旧规则失效。上一小节的模板比例示例属于第二种：变化没有制造新的组内错误，却增加了困难样本的权重。这些情形应与错误标注、测量异常等训练数据问题分开调查，因为它们所需的补救信息不同。

模型优化训练损失时，并不会自动区分稳定依据与偶然相关。若训练集中的垃圾邮件常含某个页脚，模型可能依赖页脚而非内容；当正常邮件也采用该页脚时，这条捷径就失效。\term{捷径学习}（shortcut learning）指模型借助在训练条件下有效、却不能支持预期迁移的规则完成任务。Geirhos等以此联系背景、纹理等依赖，强调应同时考察数据与评价条件\scite{geirhos2020shortcut}。

捷径不能简单等同于容易提取的特征。页脚在某些固定业务中可能确实是有效依据，问题在于使用要求是否包含页脚改变而邮件性质不变的情形。也不能把高维特征全部视为捷径、把人工可读特征全部视为稳定依据：人工规则同样可能依赖过时的相关关系。判断需要观察改变相关关系后的表现，而不是仅凭特征的外观。

一个有用的诊断是构造内容与模板不一致的样本。假设训练集中90\%的垃圾邮件使用模板甲，90\%的正常邮件使用模板乙。只读模板便能取得90\%的准确率；如果按原比例随机划分测试集，这条规则还会在测试中得到支持。将两类模板均衡配到两种邮件上后，模板规则的准确率降为50\%。这是人为构造的对照，不代表所有模型都会选择模板，却说明随机留出为何不能自动检验所需的不变性。

页脚捷径解释模型怎样出错；训练数据缺少足以打破页脚与正文关联的情形，则解释为什么训练时没有发现这条规则有问题。增加训练环境的多样性可以暴露不稳定关系，但多样性也需要与目标变化相关。如果新增邮件仍保留原来的页脚关联，即使数量很大，也没有补上页脚关系改变的情形。若目标环境的标签规律在现有数据中完全无法区分，复杂模型、正则化或无标签适应都不能凭空补上信息。定理~\ref{thm:transfer-nonidentifiability}中的不可识别反例正是这种边界。

如果任务中的真实标签关系已经改变，继续保持原输出也不再是正确目标。例如同一发件主体被接管后，系统需要根据新信息改变判断。此时应问能否及时发现并学到新规律，而不是要求对变化一律不敏感。鲁棒性要求保持的是任务表现，只有在标签保持的变化下，才能进一步把它表达成输出不变。

\section{鲁棒性的检验}
\label{sec:robustness-evaluation}

鲁棒性检验与失效处置需要区分三类结论：规定变化发生后，任务风险有多大；在一个明确邻域内，能否排除所有预测翻转；对未覆盖的输入或变化，系统能否及时识别并限制损失。前两类评价模型鲁棒性，分别依靠带标签测试和认证条件；第三类评价系统的变化应对能力，检验监测、拒绝和回退。三者相互补充：报警不代表模型必然失效，未报警也不构成稳定性证明。

\subsection{变化条件下的风险检验}
\label{sec:robustness-risk-test}

检验应从需要保持的性质出发。测量扰动测试在尽量保留任务含义的条件下改变噪声、分辨率或缺失比例；跨环境测试使用不同时间、来源或设备的数据；场景切片则检查可能被总体平均掩盖的输入。对每一种变化，都应报告强度、标签是否重新核验以及覆盖范围。

同一样本变化前后的配对比较可以减少样本组成差异，但不能替代真实新环境评价。人工构造的变化适合隔离因素，实际环境更能反映因素之间的耦合。若测试数据已经用于选择增强策略，还需要留出新的环境或扰动组合，检查是否只是适应了这组测试。

ImageNet-C把常见图像损坏组织为噪声、模糊、天气和数字处理等类型，并设置不同强度；ImageNet-P进一步考察连续小变化引起的预测不稳定。两者将“在变化下表现如何”变为可重复比较的问题，但所检验的是规定的损坏和扰动，不能据此声称对所有自然变化或最坏扰动都鲁棒\scite{hendrycks2019corruptions}。

这类基准中的归一化指标也需要读懂。例如，以参考分类器在同种损坏上的累计错误率作为分母，再对不同损坏求平均，可以减少损坏类型本身难度的影响。但得到的是相对参考模型的表现，并不是部署错误率。若实际环境以压缩损坏为主，基准中对天气和噪声的等权平均就未必对应实际收益。报告应保留按类型、强度分解的曲线，并同时给出未损坏数据的表现。

真实环境变化还包括数据来源、时间和对象群体之间的差异。WILDS汇集跨医院、相机位置、时间和地理区域等分布变化任务，并为各任务规定相应的划分和评价。它区分了两个问题：从同一来源随机留出的新样本上表现怎样，以及换到未参与训练的来源后表现怎样\scite{koh2021wilds}。

% 为完整文献边注预留高度，避免它被前面的长作者列表挤出纸边。
\Needspace{45mm}
变化还可能使不确定性表达先于准确率失效。Ovadia等将多种不确定性方法放到分布变化下比较，同时观察准确性、概率评分和校准指标怎样随变化强度改变。这类评价能够区分“预测错误增加”和“模型仍以过高把握给出错误判断”，但结论仍限于所构造的变化条件\scite{ovadia2019uncertainty}。

将邮件按月份留出也有类似作用，但要防止近重复邮件或同一会话跨越训练和测试。若目标是适应新机构，则应按机构而非逐封邮件划分；若目标是适应未来月份，则应保留时间顺序。数据怎样划分，应当取决于实际使用时会遇到哪种新情况。跨环境检验还应记录调参能够访问哪些环境，因为使用目标环境标签挑出的模型，与完全未见目标环境的模型，不能直接当作同一设定下的方法比较。

报告平均风险时，还应报告与使用相关的群体风险、样本数和不确定性。样本较少的群体通常具有更宽的统计区间，不能因为点估计较低便断定可靠。示例~\ref{ex:robustness-mixture}提示，变化环境中的期望风险还依赖组比例，因此应把组内表现与组比例变化分开。

对每个测试环境$e$，用未参与训练与调参的$n_e$个带标签样本估计
\begin{equation}
 \begin{aligned}
 \widehat R_e(f)&=\frac{1}{n_e}\sum_{i=1}^{n_e}
       \ell(f(\bm{x}_{e,i}),y_{e,i}),\\
 \widehat\Delta_e&=\widehat R_e(f)-\widehat R_0(f),
 \end{aligned}
 \label{eq:robustness-environment-test}
\end{equation}
其中$0$表示参考环境。报告绝对风险$\widehat R_e$与增加量$\widehat\Delta_e$，再按变化强度和关键群体展开，才能区分“原本就差”和“变化造成退化”。若预先约定风险容限为$\rho$，应把相应风险置信上界与$\rho$比较，而非仅看点估计是否低于它。对同一原样本生成多种扰动时，不能把相关的变体当成相互独立的样本计算区间；抽样与不确定性分析应保留原样本这一单位。

比较改进前后的模型时，应同时检验变化环境、原环境和关键群体中的表现。更新可能改善新模板却遗忘旧模板，也可能只改善检测到的变化而损害其他输入。评价数据与停止规则需要预先明确，否则持续观察并选择最好时刻同样会产生选择偏差。

没有足够证据支持更新时，保留原模型、缩小自动处理范围或回退都是控制损害的选择。它们不等于提高模型本身的鲁棒性，而是降低不确定更新的使用风险。模型能力与系统应对措施应分别报告。

部署评价可以保留冻结模型作为对照，按时间窗口比较原模型和适应模型的预测，并在标签返回后核算风险差异。若需要边运行边挑选更新时刻，置信结论必须考虑重复检验；若标签仅来自被转交人工的邮件，还要处理选择性观测，不能直接将该子集的错误率当作全量错误率。鲁棒性改进由此回到前章的证据问题：有效的变化目标、独立的检验和可解释的保证条件缺一不可。

\subsection{邻域内稳定性的认证}
\label{sec:robustness-certification}

测试集上的风险估计仍只涉及已生成的样本。若要求“半径内每一个允许输入都不翻转”，则检验对象变成一个带全称量词的命题，不能靠增加几个随机变体来替代证明。

对式~\eqref{eq:robustness-perturbation-risk}，有限搜索只提供最坏损失的下界：找到一次失败便能证实脆弱性，没有找到却不能证明不存在。若能够证明函数的变化上界，才可能获得认证。例如，二分类实值分数$f$在$U(\bm{x})$内满足
$|f(\bm{x}')-f(\bm{x})|\leq L\lVert\bm{x}'-\bm{x}\rVert_2$，
若以$\bm{x}$为中心的半径$r$球包含于$U(\bm{x})$，且$f(\bm{x})>Lr$，则球内分数保持为正。这只认证相应范数范围内的符号不变；若原预测就错了，认证不变也不会让它正确。

图~\ref{fig:trust-certified-region}把这一证明条件放到二维输入平面。认证球内的每一点都受同一个变化上界保护；球外同时存在预测不变的点与反例，因而没有证书不能被解释为已经失败。

\begin{figure}[!htb]
  \centering
  % f(x)=x1+.2x2，x=(.8,0)，L=sqrt(1.04)，认证r=.7。
\begin{tikzpicture}[figure picture,foundation picture]
 \path[use as bounding box] (0,0) rectangle (110,73);
 \begin{scope}[xshift=10mm]
 \fill[FigureBlueFill] (29,7)--(74,7)--(74,67)--(17,67)--cycle;
 \draw[draw=FigureStroke,line width=1.1pt] (29,7)--(17,67);
 \draw[draw=FigureYellow,fill=FigureYellowFill,line width=1pt] (43,37) circle (17.5);
 % 坐标轴置于区域与认证球的填充之上，保持整条轴可追踪。
 \draw[foundation axis] (4,37)--(78,37);
 \node[anchor=north] at (78,33) {$x_1$};
 \draw[foundation axis] (23,5)--(23,70) node[above] {$x_2$};
 \fill[FigureInk] (43,37) circle (1.4);
 \node[anchor=south west] at (44,39) {$\bm x=(0.8,0)$};
 \draw[draw=FigureMuted,line width=.6pt] (43,37)--(43,54.5);
 \node[anchor=west,inner sep=1.5pt] at (44,49) {$r=0.7$};
 \draw[draw=FigureBlue,fill=white,line width=1pt] (66,62) circle (1.5);
 \node[anchor=south] at (66,65) {$f>0$};
 \draw[draw=FigureRed,line width=1pt] (9,45.5)--(12,48.5) (9,48.5)--(12,45.5);
 \node[anchor=south] at (10.5,52.5) {$f<0$};
 \node[anchor=north] at (29,5) {$f=0$};
 \end{scope}
\end{tikzpicture}%

  \caption{Lipschitz条件支持的认证区域。构造$f(\bm x)=x_1+0.2x_2$，中心为$(0.8,0)$，可取$L=\sqrt{1.04}$。半径$r=0.7$满足$0.8>Lr$，所以闭球内分数始终为正。深色直线是零分数边界，认证球外的空心点仍为正，叉号则为负；球外不属于此证书的结论范围。认证保持的是预测符号，不保证原预测标签正确。}
  \label{fig:trust-certified-region}
\end{figure}

局部认证也可以针对一个经过专门构造的预测规则进行。此时必须把认证对象写清，不能将新规则的证书直接套给原模型。\term{随机平滑}（randomized smoothing）让基础分类器$h$在输入加入随机噪声后投票，以投票概率最大的类别作为新分类器的输出：
\begin{equation}
 g(\bm{x})=\arg\max_k
   \mathbb P_{\bm{\varepsilon}\sim
   \mathcal N(\bm{0},\sigma^2\bm I)}
   \{h(\bm{x}+\bm{\varepsilon})=k\}.
 \label{eq:robustness-smoothing}
\end{equation}
这里$\bm{x}\in\mathbb R^d$，$\bm I$为$d$阶单位矩阵，$\sigma>0$是噪声标准差。Cohen等建立了高斯平滑对应的$\ell_2$认证半径，关键在于平移高斯分布时，同一分类区域的概率不能任意快速改变\scite{cohen2019smoothing}。

具体地，若类别$a$的投票概率至少为$\underline p_a$，每个其他类别的概率至多为$\overline p_b$，且$\underline p_a>\overline p_b$，那么在
\begin{equation}
 \lVert\bm{\delta}\rVert_2<r_{\mathrm{cert}},
 \qquad
 r_{\mathrm{cert}}
 =\frac{\sigma}{2}
 \left[\Phi^{-1}(\underline p_a)
       -\Phi^{-1}(\overline p_b)\right]
 \label{eq:robustness-smoothing-radius}
\end{equation}
内，平滑分类器满足$g(\bm{x}+\bm{\delta})=a$。$\Phi$为标准正态分布函数。认证对象是$g$，不是任意一次加噪得到的$h$的输出，也不是未经平滑的原模型。

实际投票概率要用蒙特卡洛采样估计，因此应使用具有规定置信度的概率界。若二分类中获得$\underline p_a=0.9$，可取$\overline p_b=0.1$；令$\sigma=0.25$，认证半径约为$0.25\times1.282=0.3205$。这些数字是假设的概率界，不能把观察到90\%投票比例直接代入充当界。选类别与估计概率可以使用独立采样，界不足以分离类别时则拒绝认证。

增大$\sigma$不一定扩大实际认证半径，因为噪声更强也可能使投票概率接近，甚至损害正确率。平滑还需要多次前向计算。检验时应报告随半径变化的认证准确率，即既预测正确、又能认证该半径的样本比例，并报告采样成本与失败概率。单个样本的失败概率只控制该样本的蒙特卡洛证书；若声称一批样本或整条半径曲线同时正确，必须在全部认证事件之间分配失败概率或采用相应的联合控制。这里的概率来自认证程序对投票概率的估计；上一章的覆盖保证涉及新样本和校准数据，二者的随机对象与量词不同。

给定需要承受的开球半径$r$，实际检验应逐个样本计算预测是否正确、认证半径是否至少为$r$，再统计两项同时成立的比例。改变$r$得到认证准确率曲线，并列报告对同一预测规则进行经验搜索所得的鲁棒准确率。搜索成功找到的错误给出反例，证书成功给出该样本邻域内的保证；搜索未成功且认证也未成功的样本，则仍处于未判定状态。单点认证不自动说明测试样本代表了未来环境，关于期望风险的推断仍需前章的抽样条件。

\subsection{变化识别与拒绝效果的检验}
\label{sec:robustness-monitor-test}

对于超出已检验范围的输入，系统还可以报警或转交人工。这不能证明分类器在变化后保持正确，却能检验系统是否识别出了需要限制使用的条件。评价时应同时关注“有没有识别到变化”和“这种识别是否减少了实际损失”。

假设邮件模型主要用中文和英文正文训练，现在收到其他语言或纯图像邮件。系统需要先判断这些输入是否超出了已经验证的范围，再决定是否交给模型自动处理。\term{分布外检测}（out-of-distribution detection）用于识别相对于指定参考分布而言陌生的输入。评价前必须规定一封邮件还是一个附件算一个输入，以及哪些来源用作范围外的测试。未知类别、已有类别的语言变化和熟悉范围内的低频困难样本，并不是同一种问题；实验中的“域外”标签只表示所选测试来源，不代表所有未知分布。

为了作出接纳或转交的决定，可以利用表示空间距离、密度相关评分或模型输出构造检测分数，再在验证数据上选择阈值。但在训练分布中出现得少的输入，未必难以预测；模型给出高置信度，也不能说明输入一定来自熟悉的环境。检测器应依据漏检与误报成本评价，而不只报告排序指标；仅凭有限参考样本也不存在识别任意未知分布的统一保证。

给定一个固定的异常分数，检验应另留正常输入与规定的未知类别或来源，并在独立验证数据上选定阈值。测试时计算正常输入被误拒的比例，以及未知输入被误接纳的比例；改变阈值可以得到检测曲线，选定阈值后则应报告两种错误及其成本。接收者操作特征曲线下面积检验的是分数排序，不直接表示部署成本。文献中的FPR95有时把域内样本定义为正类，表示保留95\%域内样本时误接纳域外样本的比例，必须核对正类与分母。

未知输入也可能被正确分类，熟悉输入也可能预测错误，因此还需取得真实任务标签，在被接纳的样本上计算风险，在全体样本上计算自动处理比例，并计入转交或拒绝的成本。比较两个检测器时，应固定处理比例或成本口径。这把检测评价接回第~\ref{chap:trust-reliability}章的选择性风险检验：分布外标签用来检查陌生程度识别，任务标签用来检查是否真正减少错误，两种证据不能互换。具体检测分数与训练方法将在第~\ref{sec:robustness-rejection}节展开。

\term{漂移检测}（drift detection）则考察一段数据流是否发生统计变化。输入变化未必导致错误率上升，而输入统计稳定也可能伴随标签关系变化。没有及时标签时，漂移信号只能触发调查或更保守的策略；它不能直接说明该采用哪一种新模型。

例如，新月份的邮件模板比例改变，可能触发特征分布检验；但若模型本来就不依赖模板，预测风险未必增加。相反，原来合法的发件主体被接管后，相似输入可能对应新的风险标签，输入边缘分布检验不一定及时发现。监测中应保留少量持续获取真实标签的渠道，将输入信号与错误率、行动损失对应起来。第~\ref{chap:trust-reliability}章的统计区间适用于明确抽样条件下的风险检验，不能直接套到任意反复查看和筛选后的数据上。

检验漂移报警时，可以将保留时间顺序的数据流分成明确的无变化段与变化段。在无变化段统计误报频率，在已知变化开始时刻的测试中统计漏报比例与报警延迟，并报告窗口大小、阈值和重置规则。若真实变化时刻未知，则只能报告可核实的统计差异和后续风险，不应把每次报警都当作真阳性。获得延迟标签后，还需对照报警前后的任务风险；若报警触发了拒绝或更新，则单独核算该措施的收益和代价。

\section{鲁棒性的增强}
\label{sec:robustness-enhancement}

检验暴露薄弱环节后，改进措施应与薄弱环节对应。标签保持扰动引起的错误，需要约束局部行为；困难群体或环境捷径，需要改变训练时看重的风险与依据；已经超出已知范围的输入，需要有依据的拒绝策略；新环境中可获得的信息，则可能支持部署后适应。这些措施作用于不同对象，可以组合使用。前三者不要求持续更新部署参数，适应方法则改变运行中的模型状态，因此还承担错误积累与遗忘的风险。

选择这些约束前，还需辨别失效是否源于训练数据污染。对异常数值或标签噪声，稳健损失和重加权可以限制单个样本的影响。例如，Huber损失在小残差处保留平方惩罚，在大残差处改为线性增长，从而避免极端残差主导拟合。但真实少数群体也可能产生较大残差，机械降低其权重会使它们更难被学到。因此需要区分数据错误与模型尚未覆盖的合法样本。

\subsection{输入变化下的行为约束}
\label{sec:robustness-local-enhancement}

数据增强把预期变化转化为训练样本。对于保留标签的变换$T$，可以最小化$\ell(f(T(\bm{x})),y)$；若只有无标签样本，也可以约束变换前后输出接近。前提是变换与任务相容，过强的一致性要求可能抹掉本来有意义的差别。

当变换难以逐项列举时，可以约束样本之间的行为。\term{混合增强}（mixup）将两个样本及其标签按相同比例插值，鼓励模型在训练样本之间平缓变化。若$\bm{e}_{y}$为类别$y$的独热标签，构造
\begin{equation}
 \begin{aligned}
  \widetilde{\bm{x}}&=\lambda\bm{x}_i+(1-\lambda)\bm{x}_j,\\
  \widetilde{\bm{y}}&=\lambda\bm{e}_{y_i}
                         +(1-\lambda)\bm{e}_{y_j},
 \end{aligned}
 \label{eq:robustness-mixup}
\end{equation}
其中$\lambda$通常从对称Beta分布抽样。用软标签交叉熵训练，相当于按$\lambda$和$1-\lambda$分担两个类别的损失，而非给插值点强行指定某一端的标签\scite{zhang2018mixup}。

例如，两个原样本分别属于不同类别，取$\lambda=0.7$，插值样本的软标签就对这两类分配$(0.7,0.3)$的权重。模型在这个插值样本上给出相同概率时，便符合该训练目标希望表达的混合关系。它增加了样本之间的函数形状约束，但线性插值未必具有真实语义：把两封离散文本按字节混合通常没有意义，在表示空间插值也需要另行验证。即便原样本标签都正确，插值标签仍然是一项建模假设。混合增强因而不是式~\eqref{eq:robustness-perturbation-risk}中标签保持邻域的直接枚举。

对于图像中的常见损坏，AugMix采用另一条路线：对同一图像执行若干随机增强链，将结果混合，并保留与原图的混合通路；训练时还可约束原图和多个增强视图的预测一致。其主要目标是让模型在多种合理变化及其组合下保持稳定，而不是在两种类别之间指定线性过渡\scite{hendrycks2020augmix}。

两种增强要求模型学到的关系不同：mixup要求两个样本的插值对应标签的混合，AugMix要求同一图像的不同增强版本保持相近预测。若把增强效果只概括为“增加数据量”，就会忽略这些假设。训练前应检查变换是否破坏小目标、颜色语义或文字方向；训练后还应在未使用的损坏类型和组合上复查，避免把对某套增强算子的熟悉当作一般鲁棒性。

在标签保持的邻域内，仅随机抽样增强可能漏掉损失较高的位置。直接优化最不利邻域损失，可以迫使模型关注这些位置，但也可能改变原数据上的决策边界。TRADES把分类错误与边界附近的不稳定联系起来，并据此构造兼顾自然准确率与局部稳定性的代理目标\scite{zhang2019trades}。

以多分类概率模型$\bm{p}_{\bm{\theta}}(\bm{x})$为例，一种常用形式为
\begin{equation}
 \begin{aligned}
 \min_{\bm{\theta}}\ \mathbb E_{(\bm{x},y)}
 \bigg[&-\log p_{\bm{\theta},y}(\bm{x})\\
 &+\beta\sup_{\bm{x}'\in U(\bm{x})}
 D_{\mathrm{KL}}\!\left(
 \bm{p}_{\bm{\theta}}(\bm{x})
 \,\middle\|\,
 \bm{p}_{\bm{\theta}}(\bm{x}')\right)\bigg],
 \end{aligned}
 \label{eq:robustness-trades}
\end{equation}
其中$\beta>0$控制对稳定性的重视程度，$D_{\mathrm{KL}}$为两个类别分布的KL散度。第一项鼓励原输入为真实类别分配较高概率，第二项寻找使预测分布偏离最大的允许变化，再惩罚这种偏离。该式是连接自然分类损失与边界不稳定性的训练代理；把它解释为鲁棒零一风险上界，还需要TRADES分析中的分类校准和替代损失条件。若去掉第一项，恒定输出也能使稳定性项为零，却无法完成分类。

该目标也解释了为何不能不加区分地扩大邻域。若两个标签不同的样本邻域重叠，要求两个邻域内都保持各自标签，会产生冲突。应先核对扰动范围的任务含义，再选择权衡系数。内层搜索还需要额外计算，并通常只能近似求解；优化了这项目标不等于获得形式认证。具体如何搜索攻击样本以及检验防御是否只是遮蔽梯度，留到第~\ref{chap:trust-security}章讨论。

图~\ref{fig:trust-adversarial-training-loops}把式~\eqref{eq:robustness-trades}的两层优化分开：内层固定参数搜索更不稳定的允许输入，外层再用原输入损失与分布偏离更新参数。内层只找到搜索过程中的候选，不能把该候选写成已经证明的全局最坏输入。

\begin{figure*}[htbp]
  \centering
  % 依据已确认模型设计重建：内层固定theta，外层更新theta。
\begin{tikzpicture}[figure picture,foundation picture,
 font=\figurefont\mdseries\fontsize{9}{11}\selectfont,
 trust module/.style={draw=FigureStroke,line width=1pt,rounded corners=1.5mm,fill=white},
 trust wire/.style={draw=FigureStroke,line width=.65pt},
 trust icon/.style={draw=FigureStroke,line width=.8pt,line cap=round,line join=round},
 trust port/.style={circle,draw=FigureStroke,fill=white,line width=.8pt,minimum size=2mm,inner sep=0pt}]
 % 简化对象图标的真实路径；调用方定义trust icon和trust wire。
\providecommand{\TrustVectorMail}[3]{%
 \begin{scope}[shift={(#1,#2)},scale=#3]
  \draw[trust icon,fill=white,rounded corners=.4mm] (0,0) rectangle (8,5.5);
  \draw[trust icon] (0,5.5)--(4,2.3)--(8,5.5) (0,0)--(2.7,2.4) (8,0)--(5.3,2.4);
 \end{scope}}
\providecommand{\TrustVectorFolder}[3]{%
 \begin{scope}[shift={(#1,#2)},scale=#3]
  \draw[trust icon,fill=FigureBlueFill,rounded corners=.4mm] (0,0)--(0,5)--(3,5)--(4,4)--(8,4)--(8,0)--cycle;
 \end{scope}}
\providecommand{\TrustVectorRecord}[3]{%
 \begin{scope}[shift={(#1,#2)},scale=#3]
  \draw[trust icon,fill=white] (0,0)--(6,0)--(6,6)--(4,8)--(0,8)--cycle;
  \draw[trust icon] (4,8)--(4,6)--(6,6);
  \draw[trust icon] (1,4.8)--(4.8,4.8) (1,3.1)--(4.8,3.1) (1,1.4)--(3.8,1.4);
 \end{scope}}
\providecommand{\TrustVectorDocs}[3]{%
 \begin{scope}[shift={(#1,#2)},scale=#3]
  \foreach \dx/\dy in {0/2,1.6/1,3.2/0}{
   \draw[trust icon,fill=white,rounded corners=.5mm] (\dx,\dy) rectangle ({\dx+6},{\dy+8});
  }
  \draw[trust icon] (4.2,6)--(8.2,6) (4.2,4.1)--(8.2,4.1) (4.2,2.2)--(7.2,2.2);
 \end{scope}}
\providecommand{\TrustVectorTrash}[3]{%
 \begin{scope}[shift={(#1,#2)},scale=#3]
  \draw[trust icon,fill=FigureRedFill] (.7,0)--(0,6)--(6,6)--(5.3,0)--cycle;
  \draw[trust icon,fill=white,rounded corners=.3mm] (-.4,6) rectangle (6.4,7);
  \draw[trust icon] (2,7)--(2,8)--(4,8)--(4,7) (1.8,1)--(1.4,5) (3,1)--(3,5) (4.2,1)--(4.6,5);
 \end{scope}}
\providecommand{\TrustVectorMesh}[3]{%
 \begin{scope}[shift={(#1,#2)},scale=#3]
  \foreach \xx/\yy/\id in {-5/-2/a,-5/2/b,0/-4/c,0/0/d,0/4/e,5/-2/f,5/2/g}{\coordinate (icon-\id) at (\xx,\yy);}
  \foreach \i in {a,b}{\foreach \j in {c,d,e}{\draw[trust wire] (icon-\i)--(icon-\j);}}
  \foreach \i in {c,d,e}{\foreach \j in {f,g}{\draw[trust wire] (icon-\i)--(icon-\j);}}
  % 层内同色、层间蓝红黄；连接保持灰色，不把整网染成单色。
  \foreach \id/\col in {a/FigureBlue,b/FigureBlue,c/FigureRed,d/FigureRed,e/FigureRed,f/FigureYellow,g/FigureYellow}{\draw[trust icon,fill=\col] (icon-\id) circle (.7);}
 \end{scope}}

 \path[use as bounding box] (0,-5) rectangle (169,79);
 \draw[trust module,fill=FigureBlueFill] (0,20) rectangle (50,63);
 \draw[draw=FigureStroke,line width=.7pt,dashed] (25,41) circle (18.5);
 \node at (25,55) {$U(\bm x)$};
 \node at (25,28) {$\max\mathrm{KL}$};
 \draw[foundation flow,line width=.8pt] (18,39)--(24,42)--(29,43)--(34,47);
 \foreach \xx/\yy in {18/39,24/42,29/43}{\draw[trust icon,fill=FigureBlueFill] (\xx,\yy) circle (1);}
 \draw[trust icon,fill=white] (32.5,45.5) rectangle (35.5,48.5);
 \node at (18,34) {$\bm x$};
 \node[anchor=west] at (36,47) {$\bm x'$};
 \draw[trust module] (67,27) rectangle (95,54);
 \node[anchor=north west,inner sep=0pt] at (69,52) {$p_\theta$};
 \TrustVectorMesh{81}{40.5}{1.8}
 \draw[foundation flow] (50,44)--(66,44);
 \draw[foundation flow] (50,34)--(66,34);
 \node at (58.5,48) {$\bm x$};
 \node at (58.5,38) {$\bm x'$};
 \foreach \yy/\lab/\split in {50.5/\bm x/137,31.5/\bm x'/127}{
  \draw[trust module] (114,{\yy-3.5}) rectangle (155,{\yy+3.5});
  \fill[FigureBlueFill] (116,{\yy-2}) rectangle (\split,{\yy+2});
  \fill[FigureRedFill] (\split,{\yy-2}) rectangle (153,{\yy+2});
  \draw[draw=white,line width=.7pt] (\split,{\yy-2})--(\split,{\yy+2});
  \node at (134.5,{\yy+8}) {$p_\theta(\lab)$};
 }
 \draw[foundation flow] (95,47) to[out=0,in=180] (114,50.5);
 \draw[foundation flow] (95,35) to[out=0,in=180] (114,31.5);
 \draw[draw=FigureStroke,line width=1.1pt] (155,50.5) to[out=0,in=90] (160,41);
 \draw[draw=FigureStroke,line width=1.1pt] (155,31.5) to[out=0,in=-90] (160,41);
 \draw[foundation flow,rounded corners=2mm] (160,41)--(164,41)--(164,25)--(140,25)--(140,22);
 \draw[trust module,fill=FigureRedFill] (124,9) rectangle (158,21);
 \node at (141,15) {$\mathrm{CE}+\beta\mathrm{KL}$};
 \draw[foundation flow,rounded corners=2mm] (141,9)--(141,3)--(81,3)--(81,26);
 \node at (112,6.5) {$\min_\theta$};
 \node at (112,-1.5) {外层：更新$\theta$};
 \draw[foundation flow,rounded corners=2mm] (81,54)--(81,71)--(25,71)--(25,63);
 \node at (53,76) {内层：固定$\theta$};
 \foreach \xx/\yy in {50/44,50/34,67/44,67/34,95/47,95/35,81/54,81/27,155/50.5,155/31.5,140/21}{\node[trust port] at (\xx,\yy) {};}
\end{tikzpicture}%

  \caption{TRADES中的内层搜索与外层更新。内层在$U(\bm x)$中近似增大原输入与扰动输入预测分布的KL散度，外层最小化原输入分类损失与加权KL项，再用更新后的参数重新搜索。箭头示意搜索轨迹，概率条只示意分布发生偏离，不来自实验。该训练流程不提供全邻域的形式证书。}
  \label{fig:trust-adversarial-training-loops}
\end{figure*}

图~\ref{fig:trust-robust-accuracy-tradeoff}用有限的教学候选说明两项准确率可能怎样变化。沿A、B、C选择时存在取舍，候选D却同时改善了B的两项表现，因此这张图不能被理解为所有任务都存在同一条不可越过的准确率边界。

\begin{figure}[htbp]
  \centering
  % 固定100个教学样本，固定标签保持扰动范围；鲁棒正确集均为自然正确集子集。
% A:自然80鲁棒70；B:88/60；C:94/45；D:90/74。不是实测或普遍边界。
\begin{tikzpicture}[figure picture,foundation picture]
 \begin{axis}[width=109mm,height=79mm,xmin=74,xmax=100,ymin=35,ymax=85,
  axis lines=left,xlabel={原输入准确率（\%）},ylabel={鲁棒准确率（\%）},
  xtick={75,80,85,90,95,100},ytick={40,50,60,70,80},clip=false]
  \addplot[draw=FigureBlue,line width=.85pt,dashed,mark=o,mark size=2.7pt,
   mark options={fill=white,line width=.75pt}] coordinates {(80,70)(88,60)(94,45)};
  \node[anchor=south east] at (axis cs:79.7,71) {A};
  \node[anchor=north east] at (axis cs:87.5,60) {B};
  \node[anchor=north west] at (axis cs:94.5,45) {C};
  \addplot[only marks,draw=FigureRed,mark=diamond*,mark size=3pt,
   mark options={fill=FigureRedFill,line width=.85pt}] coordinates {(90,74)};
  \node[anchor=south west] at (axis cs:90.5,74) {D};
  \draw[foundation flow,draw=FigureRed] (axis cs:88.3,62)--(axis cs:89.7,72);
 \end{axis}
\end{tikzpicture}

  \caption{原输入准确率与鲁棒准确率的一种可能取舍。固定100个构造样本及同一标签保持扰动范围，A、B、C的两项正确数分别为80/70、88/60、94/45；D为90/74，同时优于B。每个候选的鲁棒正确样本都是原输入正确样本的子集。虚线只连接三个教学候选，不是实测曲线或普遍最优边界；评价仍需固定扰动定义与搜索条件。}
  \label{fig:trust-robust-accuracy-tradeoff}
\end{figure}

以上方法改变的是模型及其训练约束，实际收益仍应回到第~\ref{sec:robustness-risk-test}节的独立变化测试。需要严格邻域结论时，另按第~\ref{sec:robustness-certification}节执行认证；训练目标变小与认证半径变大不是同一件事。

\subsection{跨环境风险与稳定依据的约束}
\label{sec:robustness-environment-enhancement}

局部增强规定同一样本附近的行为，却未必解决困难群体被平均损失忽视、或不同环境中预测依据失效的问题。例如，少见模板的错误可能被大量常见模板掩盖，而页脚捷径可能在换环境后失效。接下来分别通过提高困难组的权重、寻找多个环境中的稳定关系，回应这两类问题。

若预期一组环境或分布都可能出现，可以用\term{分布鲁棒优化}（distributionally robust optimization, DRO）直接优化最不利分布。一般定义及其与风险决策的关系见
\BookSectionRef{sec:decision-risk-criteria-robustness}{第一卷“策略的价值与优化”章的“策略的轨迹价值与风险”一节}；
在当前监督学习问题中，这个目标写为：
\begin{equation}
 \min_f\ \sup_{\mathcal Q\in\mathfrak Q}
   \mathbb E_{\mathcal Q}\ell(f(\bm{x}),y).
 \label{eq:robustness-dro}
\end{equation}
这里$\mathfrak Q$是允许分布组成的集合。写下这个目标，还不能证明模型在实际使用中达到要求。实际使用时还须
从训练数据构造$\widehat{\mathfrak Q}(S)$，用独立数据或预先规定的统计程序选择距离、
矩约束、半径或组集合，并证明真实部署分布以声明概率落入该集合；随后还要同时控制
经验到总体的误差与内层近似优化误差。集合过小会漏掉实际变化，过大则可能要求模型抵抗
没有任务意义的情形。损失的可积性以及上确界是否可达或怎样近似，也属于保证条件。

若允许分布只是若干组分布的任意混合，式~\eqref{eq:robustness-dro}化为最小化最大组风险。它能防止大群体在平均目标中压过小群体，但经验最坏组风险很小仍不保证新样本上的组风险小。Sagawa等的研究强调，过参数化模型可以同时拟合各组训练数据，最坏组泛化仍需要正则化与验证\scite{sagawa2020groupdro}。

这个化简可以直接看出。设第$g$组的分布为$\mathcal P_g$，组比例$q_g\geq0$且$\sum_gq_g=1$，则混合分布风险为$\sum_gq_gR_g(f)$。允许任意组比例时，最不利比例会把全部权重放到风险最大的一组，因此
\begin{equation}
 \sup_{\bm{q}\in\Delta_G}\sum_{g=1}^{G}q_gR_g(f)
       =\max_{1\leq g\leq G}R_g(f),
 \label{eq:robustness-group-dro}
\end{equation}
其中$\Delta_G=\{\bm q\in\mathbb R^G:q_g\geq0,\ \sum_gq_g=1\}$是
$\mathbb R^G$中具有$G$个概率分量、仿射维数为$G-1$的概率单纯形。
示例~\ref{ex:robustness-mixture}中，普通平均目标受99\%常见模板支配，最坏组目标却会把50\%的少见模板错误率当作需要优先改善的部分。

计算时可以交替更新模型和组权重：高损失组增加权重，模型随后更重视这些组。若只允许组比例在当前比例附近变化，也可以缩小权重集合，减轻对极端混合的保守性。不过，分组太粗会掩盖组内捷径，分组太细又会产生估计不稳定的小组。组标签本身也可能不可得；仅按高损失自动分组，可能把标注错误与合法困难群体混在一起。此时需要数据诊断，不能认为优化目标会自动辨认二者。

最坏组风险与公平差异也不能互换。把所有组的损失都降到较低水平，仍可能留下机会分配或行动成本上的不合理差异；反过来，两组错误率相同，也可能意味着两组都表现很差。第~\ref{chap:trust-fairness}章将先确定需要保护的分配原则，再决定哪些组间比较有意义。本节的组DRO只针对已经给定的损失与组分布提供优化目标。

跨环境稳定关系学习进一步寻找多个环境中共同有效的依据。这个方向依赖训练环境足以打破不稳定相关。关于不变性要求何时能够支持目标域保证，可回到第~\ref{sec:domain-generalization}节；单独观察训练环境上的一致性不足以证明任意新环境有效。

不变风险最小化所寻求的是：在学习到的表示上，同一个预测规则能够同时适用于多个训练环境。它与组DRO的关注点不同；组DRO限制最坏组损失，不要求各环境的最优预测规则相同。若所有训练环境中模板与标签都同向相关，不变性目标仍可能接受模板捷径；若不同环境对目标的定义本来就不同，强求同一规则还可能删除有用信息。

方法比较因而必须包含认真调优的经验风险最小化基线。DomainBed在统一的训练和模型选择条件下比较多种域泛化方法，发现不少方法未能稳定超过强ERM基线，说明算法名称不能代替公平的训练与验证流程\scite{gulrajani2021}。

这里并不能推出ERM在所有变化下最优。它说明，在声称某种鲁棒目标有效以前，应固定架构、训练预算、数据增强和可用验证环境，检查收益究竟来自目标函数还是其他选择。第~\ref{chap:transfer-learning-theory}章给出的条件解释可识别性，本节的受控比较则检查具体算法是否把可用信息转化成实际收益。

\subsection{超出适用范围时的拒绝策略}
\label{sec:robustness-rejection}

当已有训练信息不足以支持所有输入时，可以为预测规则增加接纳条件。拒绝策略的目标是缩小缺乏证据的自动处理范围；它增强系统对变化的应对能力，并不让被拒绝输入上的原分类器自动变得正确。其代价是处理比例下降或人工负担增加，应与模型性能改善分别报告。

一个无需另训模型的基线是\term{最大类别概率}（maximum softmax probability, MSP）：用最大的预测概率作为熟悉程度分数，过低时拒绝分类。Hendrycks与Gimpel将它作为错误样本和分布外样本检测的基本比较对象\scite{hendrycks2017msp}。

这个分数只比较给定类别之间的相对支持。假设两个输入的类别打分分别是$(10,9)$和$(1,0)$，经过softmax后得到相同的最大概率，约为0.731。类别间差距相同，并不意味着两个输入都得到同样强的整体支持。这个构造也说明，上一章对类别概率作校准，不能单独解决未知类别的识别：检测需要另行规定未知的含义与测试分布。

为保留归一化前打分的信息，可以使用\term{能量分数}（energy score）。设分类器有$K$个类别，对类别$k$的实值打分为$z_k(\bm{x})$，温度$T>0$，定义
\begin{equation}
 E_T(\bm{x})=-T\log\sum_{k=1}^{K}
       \exp\!\left(\frac{z_k(\bm{x})}{T}\right).
 \label{eq:robustness-energy}
\end{equation}
通常将较高能量作为分布外信号。Liu等研究了直接用该分数检测，以及训练时拉开域内、域外能量间隔的方案\scite{liu2020energy}。

对上面的两个打分向量，取$T=1$，能量分别约为$-10.313$与$-1.313$，因而能够将它们区分。但“能够区分”不等于“区分正确”。给全部类别打分增加同一个依赖输入的常数，不改变分类概率，却改变能量；纯分类目标本身并不能唯一确定输入密度。能量阈值仍须在实际模型和验证数据上检查，不能把未经生成式拟合的能量直接宣称为真实数据的负对数密度。

如果能够取得一批辅助异常数据，还可以主动教模型区分已知任务与任务外输入。\term{异常暴露}（outlier exposure, OE）在正常分类损失外，加入对辅助异常样本的约束；分类任务中的一种实现要求这些样本的预测接近均匀分布，从而避免对它们随意给出尖锐类别判断\scite{hendrycks2019oe}。

辅助异常数据不是最终未知世界的完整代表。若训练和测试使用相同来源的异常样本，检测器可能学会来源线索而非可迁移的拒绝规则。因此应另留未参与训练的异常类别和来源，并报告相近未知类别与明显不同输入上的差异。其拒绝效果需按第~\ref{sec:robustness-monitor-test}节的独立数据和固定阈值检验。

\subsection{利用新环境信息时的风险与验收}
\label{sec:robustness-adaptation}

当新环境数据可以获得时，系统可以重新估计统计量、更新部分参数，或利用模型产生的标签调整预测规则。这类\term{测试时适应}（test-time adaptation）改变了部署中的模型状态，具体算法与训练过程留待
\BookPartRef{part:knowledge-evolution-and-transfer}{第四卷“范式”的“知识的演进与迁移”部分}
讨论；这里关注更新会怎样改变可信主张，以及需要什么证据才能继续使用模型。

没有真实标签时，适应通常依赖预测置信度、一致性或其他代理信号。代理目标改善不等于任务风险下降：错误预测可能被进一步强化，始终输出同一类别也可能得到看似稳定的无监督指标。限制更新参数、步长或样本范围可以缩小影响，却不能提供正确方向。若能取得少量真实标签，应优先用它们核验更新后的任务表现，而不是只报告代理目标。

数据流的组织也属于被验收的程序。批次大小、样本顺序、每次预测发生在更新前还是更新后、何时重置状态，都会改变后续输入所面对的模型。把测试数据随机打乱后离线更新，不能代表具有时间相关性的真实服务；只评价参与更新的高置信样本，又会遗漏被排除样本和困难群体上的风险。

适应程序应与冻结模型在同一时间顺序、输入流和资源预算下比较。验收至少同时报告原环境与新环境风险、变化轨迹、关键群体表现、更新失败后的回退效果，以及更新和保存状态的成本。模型在前一段时间的平均表现提高，不能掩盖某个短窗口中的严重退化；新环境改善也不能由旧能力遗忘来换取而不作说明。

若当前标签不足以判断更新方向，保持参考模型、缩小自动处理范围或回退，可能比继续优化代理目标更稳妥。这些措施控制的是不确定更新造成的损害，并不表示模型已经适应新环境。只有在模型、更新规则、状态与验收条件共同固定后，才能把评价结果归于这套适应程序。

\section{本章小结}
\label{sec:robustness-summary}

模型鲁棒性要求模型在规定变化下保持任务表现；系统的变化应对能力则在证据不足或模型
失效时限制损害。局部扰动、群体比例变化和标签规律变化描述评价条件；邻域边界、原有
群体误差和捷径依赖则解释这些条件怎样暴露脆弱性。两层问题分开，才能选择有意义的测试
与补救措施。常见损坏基准隔离变化因素，跨环境划分考察现实迁移，局部认证在明确邻域中
排除预测翻转；检测、拒绝和回退评价系统能否识别陌生输入并止损，却不证明原模型已经
变得鲁棒。这些证据各有范围，不能由一种测试成绩替代全部判断。

训练阶段增强与部署后适应共同回答如何提高变化条件下的有效性，但它们利用的信息和承担的风险不同。增强约束哪些输入应当具有相近行为，最坏组优化避免平均目标掩盖困难群体，局部稳定性训练处理邻域中的边界，而持续适应尝试利用新环境信息。它们都需要与认真调优的基线比较，并检查原环境、变化环境和关键群体中的得失。对于证据不足的更新，应同时考虑限制使用与回退。

开篇邮件分类器的误拦，至此可以先确定测量、群体比例或标签关系发生了什么变化，再调查邻域不稳定、既有群体误差或模板捷径如何导致损失，而无需把所有问题都归因于“数据变了”。当这些变化使错误集中在特定人群时，还需要判断不同人承担的损失是否合理。下一章将讨论公平性：在既定分配原则下，怎样比较机会、错误与负担，并改变造成不合理差异的环节。至于模型具体依赖了哪些因素，第~\ref{chap:trust-interpretability}章将进一步提供解释与核验的方法。

\paragraph{延伸阅读}

Hendrycks与Dietterich关于常见损坏的基准、Koh等人的WILDS任务集，适合对照受控扰动与真实来源变化怎样形成不同评价证据。阅读时应保留变化类型、数据划分和模型选择条件，避免把某个基准上的平均结果推广到所有环境\scite{hendrycks2019corruptions,koh2021wilds}。

若希望继续研究改进方法，可从TRADES与组DRO比较局部稳定性和最坏组风险的不同目标；测试时适应的具体算法、状态更新和实践选择留到
\BookPartRef{part:knowledge-evolution-and-transfer}{第四卷“范式”的“知识的演进与迁移”部分}
展开\scite{zhang2019trades,sagawa2020groupdro}。
